% ============================================================
% Capítulo 1 — Introducción
% ============================================================
\chapter{Introducción a la reidentificación de personas}
\label{cap:introduccion}

Este capítulo sienta las bases del problema que el modelo intenta
resolver. Toda la información proviene del paper traducido al español
en \file{docs/paper.md}.

\section{¿Qué es la reidentificación de personas?}

La \emph{reidentificación de personas} (en inglés \emph{person
re-identification}, abreviada \reid{}) es la tarea de reconocer a la
\emph{misma} persona en distintas fotografías capturadas por cámaras
diferentes -- o por la misma cámara en distintos momentos. La consulta
de búsqueda puede ser:

\begin{itemize}
  \item una \textbf{imagen} (\emph{image-to-image re-ID}),
  \item una \textbf{secuencia de vídeo},
  \item una \textbf{descripción textual} (\emph{text-based re-ID}).
\end{itemize}

A diferencia del reconocimiento facial, la \reid{} no depende de
características faciales -- que pueden ser de baja calidad o estar
restringidas por privacidad -- sino de atributos visibles como ropa,
color, tipo de cuerpo o accesorios \cite{agyeman2024decentralized}.
Esto la hace especialmente útil en entornos concurridos donde es
difícil obtener caras nítidas.

\section{Reidentificación basada en texto}

La \reid{} basada en texto (\emph{text-based person re-identification})
consiste en recuperar la imagen de una persona a partir de una
descripción en lenguaje natural. Por ejemplo, ante la consulta:

\begin{quote}
\textit{``Una mujer lleva una camiseta negra con letras rosas. Viste
pantalones cortos color caqui y zapatos blancos.''}
\end{quote}

\noindent el sistema debe devolver, de una galería de miles de
imágenes, las que correspondan a esa persona.

Esta modalidad tiene dos ventajas prácticas frente a la \reid{}
basada en imagen:

\begin{enumerate}
  \item Las consultas son más \textbf{accesibles}: cualquiera puede
        describir a una persona, pero no siempre se tiene una foto
        reciente.
  \item Resuelve \textbf{problemas de privacidad}: en muchas
        jurisdicciones capturar y almacenar caras está restringido,
        mientras que las descripciones textuales son menos invasivas.
\end{enumerate}

\section{Desafíos de los sistemas existentes}

Los algoritmos \emph{state-of-the-art} de \reid{} por texto suelen
ser modelos enormes (cientos de millones de parámetros) que
priorizan precisión sobre eficiencia. Esto trae tres problemas en
aplicaciones reales:

\begin{itemize}
  \item \textbf{Requisitos computacionales altos}: no se pueden
        ejecutar en dispositivos embebidos como una cámara
        inteligente.
  \item \textbf{Dependencia de arquitecturas centralizadas}: todas
        las cámaras envían vídeo a un servidor central que ejecuta
        la identificación. Esto genera cuellos de botella,
        problemas de escalabilidad y consumo elevado de ancho de
        banda.
  \item \textbf{Imágenes pre-recortadas}: la mayoría de los
        algoritmos asumen que la imagen de entrada contiene una
        sola persona ya recortada, lo cual es irreal en escenas
        concurridas.
\end{itemize}

\section{La solución propuesta: U-TextReIDNet}

El paper de Agyeman \emph{et al.}~\cite{agyeman2024decentralized}
propone \textbf{U-TextReIDNet} (Unified Text Re-Identification
Network), un modelo:

\begin{itemize}
  \item \textbf{Preciso}: alcanza \textbf{54,02\%} Top-1 en
        CUHK-PEDES y \textbf{38,45\%} Top-1 en RSTPReid.
  \item \textbf{Ligero}: solo \textbf{38,77 millones} de
        parámetros totales.
  \item \textbf{Eficiente}: corre en tiempo real en un
        \emph{NVIDIA Jetson Nano}.
  \item \textbf{Unificado}: integra la detección de personas con
        la identificación, eliminando la necesidad de pre-recortar.
\end{itemize}

A diferencia de sus predecesores \emph{MHParsNet} (detección) y
\emph{TextReIDNet} (identificación), que eran módulos separados,
U-TextReIDNet los une en un solo modelo end-to-end.

\section{El dataset CUHK-PEDES}

Este proyecto entrena sobre \textbf{CUHK-PEDES}
\cite{li2017personsearch}, uno de los benchmarks estándar para
\reid{} por texto. Estadísticas según el paper:

\begin{table}[H]
\centering
\caption{Estadísticas del dataset CUHK-PEDES original.}
\label{tab:cuhkpedes-stats}
\begin{tabular}{lrrr}
\toprule
\textbf{Split} & \textbf{Imágenes} & \textbf{Descripciones} & \textbf{Personas} \\
\midrule
Train & 34\,054 & 68\,126 & 11\,003 \\
Val   &  3\,078 &  6\,158 &  1\,000 \\
Test  &  3\,074 &  6\,156 &  1\,000 \\
\midrule
\textbf{Total} & \textbf{40\,206} & \textbf{80\,440} & \textbf{13\,003} \\
\bottomrule
\end{tabular}
\end{table}

\noindent Cada imagen tiene asociadas \emph{dos} descripciones
textuales escritas por anotadores humanos. En este proyecto se usa
una versión alternativa descargada de \emph{HuggingFace}
(\texttt{PeterPanTheGenius/CUHK-PEDES}) que tiene una sola
descripción por imagen pero está disponible sin necesidad de
solicitar acceso por email. Los detalles están en el
\cref{cap:dataset}.

\section{Lo que este libro no cubre}

U-TextReIDNet completo (paper) incluye dos componentes:

\begin{itemize}
  \item \textbf{R-MHParsNet}: detector de múltiples personas
        (\emph{FPN} + EfficientNet-B0 + cabezas de segmentación).
        Tiene \textbf{6,49M} parámetros y NO está implementado en
        este proyecto.
  \item \textbf{TextReIDNet}: el modelo de identificación
        imagen-texto (EfficientNet-B0 + BERT + BiGRU). Éste es el
        núcleo del proyecto \texttt{textreid-train}.
\end{itemize}

Este libro se centra exclusivamente en \textbf{TextReIDNet} y su
entrenamiento. Para la parte descentralizada (Jetson Nano, ZMQ,
MQTT, Flask) se da contexto en el \cref{cap:sistema-descentralizado}
pero sin entrar en código, ya que el repositorio no la incluye.
